\chapter[Analyse des Besoins \& Conception]{Analyse des Besoins \& Conception du Système}
\label{chap:conception}

\section{Analyse des Besoins}
La conception de CyberScore TN repose sur une étude rigoureuse des exigences fonctionnelles et opérationnelles formulées par les officiers de sécurité de l'ANCS et des ministères.

\subsection{Besoins Fonctionnels}
Le système couvre seize exigences fonctionnelles prioritaires réparties selon leur niveau de criticité opérationnelle :

\begin{table}[!htbp]
\centering
\caption{Matrice des 16 exigences fonctionnelles de CyberScore TN}
\label{tab:besoins_fonctionnels}
\footnotesize
\renewcommand{\arraystretch}{1.08}
\begin{tabularx}{\textwidth}{l X c p{3.8cm}}
\toprule
\textbf{Code} & \textbf{Intitulé du Besoin Fonctionnel} & \textbf{Priorité} & \textbf{Module Réalisateur} \\
\midrule
\textbf{BF01} & Audit passif automatique des 120+ domaines gouvernementaux & Haute & Moteur d'Orchestration \\
\textbf{BF02} & Calcul mathématique du score sur 100 avec notation duale (A--F, $\star\star\star$) & Haute & Score Engine \\
\textbf{BF03} & Mappage automatisé vers OWASP Top 10, NIST CSF v2 et ISO 27001 & Haute & Framework Mapper \\
\textbf{BF04} & Visualisation télémétrique au sein du Command Center National & Haute & Next.js Dashboard \\
\textbf{BF05} & Cartographie vectorielle SVG des 24 gouvernorats tunisiens & Haute & TunisiaMap SVG \\
\textbf{BF06} & Génération et export d'attestations d'audit PDF vectorielles 7 pages & Haute & React-PDF Engine \\
\textbf{BF07} & Dialogue souverain avec l'analyste IA 3D Amina (streaming SSE) & Haute & Ollama / Qwen2.5-Coder \\
\textbf{BF08} & Énumération passive de la surface d'attaque sous-domaines (EASM) & Moyenne & Subfinder / DNS \\
\textbf{BF09} & Surveillance CTI des cybermenaces mondiales (War Room Canvas 60 FPS) & Moyenne & Threat Intel Engine \\
\textbf{BF10} & Identification des technologies serveurs et des CMS obsolètes & Moyenne & Tech Detector \\
\textbf{BF11} & Détection des faiblesses Supply Chain et des scripts CDN sans SRI & Moyenne & Supply Chain Scanner \\
\textbf{BF12} & Modélisation prédictive de la dérive de score à +30 jours & Moyenne & Prediction Analytics \\
\textbf{BF13} & Émulation professionnelle des RSSI par attribution d'XP et 12 badges & Basse & Gamification Hub \\
\textbf{BF14} & Génération du Sceau Public de Confiance avec QR code anti-fraude & Basse & Trust Seal Generator \\
\textbf{BF15} & Observatoire comparatif et benchmarking géostratégique (MENA \& Monde) & Basse & Benchmark Engine \\
\textbf{BF16} & Contrôle d'accès strict RBAC avec interception 403 et Audit Trail & Haute & Security Middleware \\
\bottomrule
\end{tabularx}
\end{table}

\subsection{Besoins Non-Fonctionnels}
Les spécifications non-fonctionnelles garantissent l'excellence technique et la viabilité étatique de la plateforme :

\begin{table}[!htbp]
\centering
\caption{Matrice des exigences non-fonctionnelles}
\label{tab:besoins_non_fonctionnels}
\small
\renewcommand{\arraystretch}{1.1}
\begin{tabularx}{\textwidth}{l p{4cm} X}
\toprule
\textbf{Code} & \textbf{Exigence} & \textbf{Spécification Technique Validée} \\
\midrule
\textbf{BNF01} & Souveraineté Numérique Totale & 100\% des données et de l'inférence IA hébergées en local ; 0 octet transmis vers des clouds tiers. \\
\textbf{BNF02} & Vélocité d'Exécution Asynchrone & Scan Quick en $\sim$2.8 s, Standard en $\sim$7.4 s et Deep en $\sim$28.6 s via parallélisation \texttt{asyncio.gather}. \\
\textbf{BNF03} & Ergonomie Trilingue \& RTL & Prise en charge native de l'arabe avec inversion symétrique CSS Logical Properties sans latence. \\
\textbf{BNF04} & Portabilité \& Conteneurisation & Déploiement reproductible orchestré sous Docker Compose articulant 11 microservices isolés. \\
\textbf{BNF05} & Résilience \& Tolérance aux Pannes & Basculement transparent vers un moteur expert déterministe certifié en cas de coupure du GPU LLM. \\
\textbf{BNF06} & Accessibilité Multi-Résolutions & Interface réactive fluide adaptée aux murs d'écrans SOC 4K comme aux terminaux portables d'astreinte. \\
\bottomrule
\end{tabularx}
\end{table}

\section{Architecture Globale à 4 Couches}
Pour répondre à ces exigences avec une modularité maximale, nous avons structuré CyberScore TN selon une \textbf{Architecture Logicielle à 4 Couches Distinctes} :

\begin{figure}[!htbp]
\centering
\begin{tcolorbox}[colback=gray!5,colframe=tekupPurple,arc=1mm,title=\textbf{Architecture Modulaire à 4 Couches de CyberScore TN}]
\centering
\textbf{\large 1. COUCHE PRÉSENTATION (FRONTEND)}\\[1mm]
Next.js 16 (App Router) $\cdot$ React 19 $\cdot$ Three.js 3D Avatar (Amina) $\cdot$ RTL Arabe $\cdot$ Canvas 60 FPS\\
$\downarrow$ \textit{\small Protocole REST (JSON) + Flux Streaming SSE (Server-Sent Events) + WebSockets}\\[2mm]
\textbf{\large 2. COUCHE SERVICE \& MOTEURS D'AUDIT (BACKEND)}\\[1mm]
FastAPI (Python 3.12) $\cdot$ AsyncIO $\cdot$ Pydantic v2 $\cdot$ 8 Modules de Scan $\cdot$ Moteur de Notation sur 100 $\cdot$ RBAC\\
$\downarrow$ \textit{\small Appels Asynchrones \& Interrogation Sockets / Modèles de Langage}\\[2mm]
\textbf{\large 3. COUCHE INTELLIGENCE ARTIFICIELLE SOUVERAINE (IA LOCALE)}\\[1mm]
Runtime Ollama Local $\cdot$ Modèle Qwen2.5-Coder:7b $\cdot$ Générateur Expert de Remédiation Déterministe\\
$\downarrow$ \textit{\small ORM SQLAlchemy 2.0 (asyncpg) $\cdot$ Cache In-Memory}\\[2mm]
\textbf{\large 4. COUCHE DONNÉES \& PERSISTANCE}\\[1mm]
PostgreSQL 16 (19 Tables Relationnelles \& JSONB) $\cdot$ Redis 7 (Cache de Scan \& Files Celery)
\end{tcolorbox}
\caption{Vue architecturale à 4 couches de la plateforme CyberScore TN}
\label{fig:archi_4_couches}
\end{figure}

\section{Modélisation UML du Système}

\subsection{Diagramme des Cas d'Utilisation (Use Case)}
Le système distingue trois profils d'utilisateurs rattachés à l'État :
\begin{itemize}
    \item \textbf{L'Administrateur Central ANCS (Super-Admin)} : Supervise l'ensemble des 120 institutions, valide ou révoque les comptes officiers, configure les seuils d'alerte nationaux et active le mode maintenance.
    \item \textbf{L'Auditeur / Officier RSSI Ministériel} : Dispose d'une habilitation restreinte à son périmètre ministériel, lance des audits de sécurité, dialogue avec l'IA Amina et exporte les attestations PDF officielles.
    \item \textbf{L'Observateur Public} : Accède au Command Center et au Benchmark international en lecture seule, sans privilège de scan ni accès aux configurations internes.
\end{itemize}

\begin{figure}[!htbp]
\centering
\begin{tcolorbox}[colback=white,colframe=tekupPurple,arc=1mm,left=4pt,right=4pt,top=4pt,bottom=4pt,title=\textbf{Matrice des Droits d'Accès par Rôle (RBAC)}]
\begin{singlespace}
\footnotesize
\begin{tabularx}{\linewidth}{X >{\centering\arraybackslash}p{2.6cm} >{\centering\arraybackslash}p{2.4cm} >{\centering\arraybackslash}p{2.6cm}}
\toprule
\textbf{Cas d'Utilisation Métier} & \textbf{Admin ANCS} & \textbf{Auditeur RSSI} & \textbf{Observateur} \\
\midrule
Command Center \& Cartographie & \checkmark & \checkmark & \checkmark \\
Benchmark International & \checkmark & \checkmark & \checkmark \\
Lancement Scan (/scan) & \checkmark & \checkmark & $\times$ \textit{(Bloqué)} \\
Assistant IA Amina & \checkmark & \checkmark & $\times$ \\
Export Rapport PDF Officiel & \checkmark & \checkmark & $\times$ \\
Gestion des Utilisateurs & \checkmark & $\times$ \textit{(HTTP 403)} & $\times$ \textit{(HTTP 403)} \\
Mode Maintenance National & \checkmark & $\times$ & $\times$ \\
\bottomrule
\end{tabularx}
\end{singlespace}
\end{tcolorbox}
\caption{Matrice d'attribution des cas d'utilisation par rôle}
\label{fig:rbac_matrix}
\end{figure}

\subsection{Diagramme de Séquence -- Cycle de Vie d'un Audit Asynchrone}
Le déroulement d'un audit de sécurité illustre la puissance de la parallélisation asynchrone :
\begin{enumerate}
    \item L'auditeur transmet l'URL cible (ex: \texttt{finances.gov.tn}) et sélectionne le mode d'audit via l'interface Next.js.
    \item L'API FastAPI réceptionne la requête et mandate l'Orchestrateur asynchrone (\texttt{orchestrator.py}).
    \item Grâce à la coroutine \texttt{asyncio.gather}, les modules d'inspection s'exécutent simultanément : poignée de main TLS (\texttt{sslyze}), inspection des en-têtes (\texttt{httpx}), résolution DNSSEC/DMARC (\texttt{dnspython}) et analyse des cookies.
    \item Les constats bruts normalisés sous forme d'objets \texttt{Finding} sont transmis au Score Engine.
    \item Le Score Engine calcule la note sur 100, attribue le grade (ex: B+) et le label d'étoiles, puis sauvegarde la transaction en base PostgreSQL 16.
    \item L'interface affiche instantanément les gauges et le radar chart, et offre l'option de consulter l'IA Amina pour la remédiation.
\end{enumerate}
